\documentclass[11pt]{article}

% ===== 基本宏包：尽量保持简单 =====
\usepackage[margin=1in]{geometry}
\usepackage{graphicx}
\usepackage{float}
\usepackage{chngcntr}
\usepackage{booktabs}
\usepackage{multirow}
\usepackage{longtable}
\usepackage{array}
\usepackage{calc}
\usepackage{amsmath, amssymb}
\setlength{\LTcapwidth}{\textwidth}
\usepackage{xurl}
\usepackage{hyperref}
\usepackage[numbers,sort&compress]{natbib}
\usepackage{setspace}
\usepackage{fontspec}
\usepackage{xeCJK}
\providecommand{\tightlist}{\setlength{\itemsep}{0pt}\setlength{\parskip}{0pt}}

% ===== 图表编号风格：全文连续编号，使用 Figure / Table =====
\renewcommand{\figurename}{Figure}
\renewcommand{\tablename}{Table}

% ===== 字体与换行策略 =====
\defaultfontfeatures{Ligatures=TeX,Scale=MatchLowercase}
\IfFontExistsTF{Noto Serif}{\setmainfont{Noto Serif}}{\IfFontExistsTF{DejaVu Serif}{\setmainfont{DejaVu Serif}}{\setmainfont{Latin Modern Roman}}}
\IfFontExistsTF{Noto Sans}{\setsansfont{Noto Sans}}{\IfFontExistsTF{DejaVu Sans}{\setsansfont{DejaVu Sans}}{\setsansfont{Latin Modern Sans}}}
\IfFontExistsTF{Noto Sans Mono}{\setmonofont{Noto Sans Mono}}{\IfFontExistsTF{DejaVu Sans Mono}{\setmonofont{DejaVu Sans Mono}}{\setmonofont{Latin Modern Mono}}}
\IfFontExistsTF{Noto Serif CJK SC}{\setCJKmainfont[AutoFakeBold=3,AutoFakeSlant=.2]{Noto Serif CJK SC}}{\IfFontExistsTF{Source Han Serif SC}{\setCJKmainfont[AutoFakeBold=3,AutoFakeSlant=.2]{Source Han Serif SC}}{\IfFontExistsTF{AR PL UMing CN}{\setCJKmainfont[AutoFakeBold=3,AutoFakeSlant=.2]{AR PL UMing CN}}{\setCJKmainfont[AutoFakeBold=3,AutoFakeSlant=.2]{WenQuanYi Zen Hei}}}}
\IfFontExistsTF{Noto Sans CJK SC}{\setCJKsansfont[AutoFakeBold=3,AutoFakeSlant=.2]{Noto Sans CJK SC}}{\IfFontExistsTF{Source Han Sans SC}{\setCJKsansfont[AutoFakeBold=3,AutoFakeSlant=.2]{Source Han Sans SC}}{\IfFontExistsTF{Droid Sans Fallback}{\setCJKsansfont[AutoFakeBold=3,AutoFakeSlant=.2]{Droid Sans Fallback}}{\setCJKsansfont[AutoFakeBold=3,AutoFakeSlant=.2]{AR PL UMing CN}}}}
\IfFontExistsTF{Noto Sans Mono CJK SC}{\setCJKmonofont[AutoFakeBold=3]{Noto Sans Mono CJK SC}}{\IfFontExistsTF{Droid Sans Fallback}{\setCJKmonofont[AutoFakeBold=3]{Droid Sans Fallback}}{\setCJKmonofont[AutoFakeBold=3]{AR PL UMing CN}}}
\setlength{\columnsep}{18pt}
\setlength{\emergencystretch}{6em}
\sloppy

% 可选：让行距略宽一些，便于审稿阅读
\setstretch{1.05}

\title{\textbf{Decathlon VOC Analyzer: An Evidence-Driven Multimodal VOC Analysis System for Aligning Product Images, Product Text, and User Reviews}}

\author{{\small
\begin{tabular}{@{}ccc@{}}
\parbox[t]{0.3067\textwidth}{\centering \textbf{Severin Ye}\textsuperscript{1}} & \parbox[t]{0.3067\textwidth}{\centering \textbf{Dokeun Lee}\textsuperscript{2}} & \parbox[t]{0.3067\textwidth}{\centering \textbf{Wushuang Liu}\textsuperscript{3}} \\[0.9em]
\parbox[t]{0.3067\textwidth}{\centering \textbf{Seowan Jung}\textsuperscript{1}} & \parbox[t]{0.3067\textwidth}{\centering \textbf{HyunJun Jung}\textsuperscript{1}} & \parbox[t]{0.3067\textwidth}{\centering \textbf{Hye Jeon}\textsuperscript{1}} \\[0.9em]
\parbox[t]{0.3067\textwidth}{} & \parbox[t]{0.3067\textwidth}{\centering \textbf{Jaesoo Kim}\textsuperscript{1*†}} & \parbox[t]{0.3067\textwidth}{} \\
\multicolumn{3}{c}{} \\[0.45em]
\multicolumn{3}{c}{\parbox{0.9200\textwidth}{\centering\small\textsuperscript{1}\textit{School of Computer Science and Engineering, Kyungpook National University}}} \\
\multicolumn{3}{c}{\parbox{0.9200\textwidth}{\centering\small\textsuperscript{2}\textit{Department of Physics, Kyungpook National University}}} \\
\multicolumn{3}{c}{\parbox{0.9200\textwidth}{\centering\small\textsuperscript{3}\textit{Department of Child Studies, Kyungpook National University}}} \\
\multicolumn{3}{c}{\small \textsuperscript{†}  Corresponding Author} \\
\end{tabular}
}}

\date{} % Nature 系列投稿一般不需要显示日期

\begin{document}
\maketitle
\nocite{lewis2020retrievalAugmentedGeneration,gao2023rarr,radford2021clip,liu2024structuredOutputBenchmark,geng2023grammarConstrainedDecoding,yao2023react,nvidia2024multimodalRagIntro,nvidia2024aiReadyKnowledgeSystems,faysse2024colpali,pontiki2014semevalAspect,qdrant2024pdfRetrievalAtScale,qdrant2024multiVectorCourse,manning2008introductionIr}

% =========================
% Abstract
% =========================
\begin{abstract}
E-commerce voice-of-customer (VOC) analysis requires understanding subjective customer feedback and grounding it in objective product-page evidence. Text-only review analysis can identify sentiment and aspects, but it cannot determine whether a user claim is supported by product copy, specifications, images, or visual details. Conversely, directly asking a large language model to summarize all reviews and product information often produces fluent but weakly auditable conclusions. This paper presents Decathlon VOC Analyzer, an evidence-driven multimodal VOC analysis system for aligning customer reviews with product text and image evidence.

The method normalizes raw product information into traceable evidence packages, including product text, product images, default image regions, and review records. It performs review filtering, rating-aware sampling, aspect-level structured extraction, and deduplication. Each extracted aspect is then converted into evidence-seeking questions with explicit text, image, or cross-modal retrieval routes. The retrieval layer performs two-stage recall and reranking over product text and image evidence. The generation layer produces structured reports containing strengths, weaknesses, controversies, evidence gaps, improvement suggestions, and claim-level evidence attribution. Key intermediate representations and evaluation records are retained to support reproducibility and error analysis.

The main contributions are threefold. First, we propose a product VOC framework based on review aspect modeling, question planning, multimodal evidence retrieval, and evidence-constrained generation. Second, we implement an artifact-first research system in which reviews, questions, evidence, reports, and evaluation records are represented by structured schemas. Third, we provide a system validation protocol that connects automated tests, run manifests, retrieval-quality metrics, and claim-attribution metrics through a common evaluation interface. The paper is positioned as a system-methodology paper focused on architecture, implementation mechanisms, and reproducible evaluation foundations.

Keywords: multimodal RAG; voice of customer; aspect-level review modeling; evidence-constrained generation; product review analysis; evidence attribution
\end{abstract}

% =========================
% Introduction
% =========================
\section{Introduction}
Customer review analysis is a core task in e-commerce intelligence. Traditional approaches usually focus on sentiment classification, topic modeling, or aspect-based sentiment analysis. They can summarize what users like or dislike, but they rarely answer whether a claim is supported, contradicted, or left unresolved by the product page itself. For product operations, knowing that users complain about size or comfort is not sufficient; analysts also need to know whether the product copy states the relevant specifications and whether product images expose the corresponding structure.

Large language models make review summarization easier, but one-shot summarization introduces a different risk. If reviews, product text, and image descriptions are simply concatenated and passed to a model, the resulting report may be fluent while remaining difficult to audit. Users cannot reliably determine which claims came from reviews, which came from product text, and which were inferred from images. Error localization also becomes difficult.

More broadly, existing product review analysis methods often remain at the level of sentiment statistics, topic clustering, aspect-sentiment extraction, or black-box summarization. These methods can provide an overall view of customer attitudes, but they rarely model the correspondence between review claims and product evidence. They also tend not to distinguish product defects, insufficient page presentation, and expectation mismatch. In addition, if review sampling is unconstrained, high-rating short comments or low-information reviews can consume the limited analysis budget, weakening the representation of improvement-oriented signals. Therefore, this paper focuses not only on summarizing product reviews, but also on prioritizing diagnostically valuable VOC signals under budget constraints and grounding those signals in verifiable product text and image evidence.

This paper studies how to explicitly align subjective customer feedback with verifiable product text and image evidence. We propose Decathlon VOC Analyzer, a staged system that constructs product-centered evidence packages, extracts review aspects, plans evidence-seeking questions, retrieves product text and image evidence, reranks candidates, and generates structured VOC reports with claim-level attribution.

The contributions are as follows. First, we propose an evidence-driven multimodal RAG framework for product VOC analysis, adding a question-planning layer between review aspects and product evidence. Second, we implement an artifact-first prototype in which evidence packages, review aspects, retrieval questions, retrieval records, reports, feedback, replay files, and manifests are all structured and persistent. Third, we provide reproducible system interfaces and evaluation hooks for both retrieval metrics and generated-claim grounding metrics.

The rest of the paper is organized as follows. Section 2 defines the task and design constraints. Section 3 reviews related work. Section 4 describes the methodology. Section 5 presents the experimental setup. Section 6 reports system validation results. Section 7 discusses implications. Sections 8 and 9 conclude and describe limitations.

% =========================
% Background and Problem Definition
% =========================
\section{Background and Problem Definition}
\subsection{Evidence Needs in Product VOC Analysis}

Product VOC analysis is not equivalent to generic review summarization. A useful product report should contain stable analytical objects such as strengths, weaknesses, controversies, applicable scenes, evidence gaps, and improvement suggestions. Each insight should ideally trace back to both customer review evidence and product-page evidence.

If a system only analyzes reviews, it answers what customers say but not whether the product page explains or supports those claims. If a system only retrieves product-page evidence, it lacks the anchor of actual customer concerns. Therefore, the key challenge is connecting subjective review needs with objective product evidence.

\subsection{Input and Output}

Given a product package \(P\) containing product text \(P_{text}\), product images \(P_{image}\), and reviews \(C_{reviews}\), the system outputs a structured report \(R\):

\[
R = F(P_{text}, P_{image}, C_{reviews})
\]

In this paper, \(F\) is not a one-step generation function. It is a staged workflow:

\[
C_{reviews} \rightarrow A_{aspects} \rightarrow I_{intents} \rightarrow Q_{questions} \rightarrow E_{retrieval} \rightarrow G_{aggregates} \rightarrow R_{attributed}
\]

where \(A_{aspects}\) denotes review aspects, \(I_{intents}\) denotes question intents, \(Q_{questions}\) denotes executable retrieval questions, \(E_{retrieval}\) denotes retrieved product evidence, \(G_{aggregates}\) denotes aspect-level aggregation, and \(R_{attributed}\) denotes the final attributed report.

\subsection{Design Constraints}

The system follows five constraints. Product evidence must be object-based and traceable. Reviews must be structured into aspects before retrieval. Retrieval queries should be questionized instead of using raw review text. Text and image evidence should be managed as separate routes. Final report claims should be traceable to review IDs, text block IDs, image IDs, or region IDs; unsupported claims should be represented as evidence gaps.

\subsection{Difference from Generic RAG QA}

Generic RAG systems usually answer a user question. In this task, the query source is a set of review aspects, and the output is a product VOC profile rather than a single answer. The system must manage multiple aspects, questions, evidence routes, report fields, and trace objects, making it closer to an evidence orchestration system than a standard QA pipeline.

% =========================
% Related Work
% =========================
\section{Related Work}
\subsection{Retrieval-Augmented Generation and Evidence Constraints}

Retrieval-augmented generation separates external evidence retrieval from parametric generation, improving traceability in knowledge-intensive tasks \citep{lewis2020retrievalAugmentedGeneration}. Retrieval-grounded revision further shows that external evidence can constrain and revise model outputs \citep{gao2023rarr}. These ideas motivate our treatment of product VOC reporting as evidence-constrained generation rather than free-form summarization.

\subsection{Multimodal Retrieval and Visual Evidence}

CLIP demonstrates that images and text can be embedded into a shared semantic space for cross-modal retrieval \citep{radford2021clip}. Engineering discussions of multimodal RAG also emphasize that real-world knowledge systems contain text, images, tables, and structured fields, requiring representation, retrieval, and fusion across modalities \citep{nvidia2024multimodalRagIntro,nvidia2024aiReadyKnowledgeSystems}. Visual document retrieval methods such as ColPali further argue for preserving original visual evidence and using finer-grained visual matching \citep{faysse2024colpali}. Our system does not implement ColPali-style late interaction, but it preserves product images as first-class evidence and creates default image regions as a lightweight interface for local visual evidence.

\subsection{Review Analysis and Aspect Modeling}

Aspect-based sentiment analysis decomposes reviews into attributes, opinions, and sentiment polarity \citep{pontiki2014semevalAspect}. In product VOC analysis, this remains a necessary first step because reviews often mix multiple experiences, scenes, and emotions. In our system, however, aspects are not the final output. They become the starting point for evidence-seeking question generation.

\subsection{Multi-Stage Retrieval and Reranking}

Practical retrieval systems often separate coarse recall from expensive reranking. Qdrant engineering materials emphasize that vector stores should manage candidates and similarity search, while rerankers and downstream logic make higher-precision decisions \citep{qdrant2024pdfRetrievalAtScale,qdrant2024multiVectorCourse}. Our system follows this separation by decoupling embedding services, index backends, candidate pooling, reranking, and final report attribution.

\subsection{Structured Output and Workflow Orchestration}

LLMs can be brittle in structured-output tasks, especially when outputs must satisfy stable field and type constraints \citep{liu2024structuredOutputBenchmark,geng2023grammarConstrainedDecoding}. ReAct-style reasoning and acting also suggests that multi-step stateful workflows are preferable to one-shot prompting for complex tool-mediated tasks \citep{yao2023react}. Our system uses structured schemas for review extraction, question generation, report generation, caching, and evaluation, and uses LangGraph to orchestrate the single-product workflow.

\subsection{Positioning}

This paper does not propose a new general-purpose embedding model or reranker. Its contribution is a system-level framework for product VOC analysis that combines review aspects, question planning, multimodal retrieval, structured reporting, and claim-level attribution into a reproducible evidence workflow.

% =========================
% Methodology
% =========================
\section{Methodology}
\subsection{Overall Framework}

Decathlon VOC Analyzer formulates product VOC analysis as an evidence-driven multi-stage reasoning process. The system first organizes product text, images, and reviews into a unified evidence space. It then extracts aspect-level customer feedback, converts each aspect into evidence-seeking questions, retrieves and reranks product text and image evidence, and finally generates structured reports with claim-level attribution. Unlike one-step LLM summarization, the framework separates what customers express, whether product evidence supports it, and how the final report grounds its claims.

The framework contains four logical layers. The product evidence layer constructs structured representations of product text, images, and local visual regions. The VOC demand layer converts reviews into aspects, opinions, sentiments, and usage scenes. The evidence alignment layer rewrites aspects into route-aware questions and retrieves evidence from the product evidence space. The report attribution layer generates structured insights and maps report claims back to review and product evidence.

\begin{figure}[H]
\centering
\includegraphics[width=0.92\textwidth]{/home/severin/Codelib/Decathlon\_VOC\_Analyzer/论文\_md形式/图片/图1.png}
\caption{Overall framework of Decathlon VOC Analyzer. The system uses review aspects and product evidence packages as intermediate representations, connects subjective VOC signals with multimodal product evidence through question planning, and performs claim-level attribution during reporting.}
\end{figure}

Table 1 summarizes the differences between the proposed system and traditional review summarization or aspect-based sentiment analysis in terms of output objective, evidence access, review handling, and traceability. The comparison emphasizes that this work does not treat aspect extraction as the final output. Instead, aspect objects serve as the intermediate entry point for evidence retrieval, report generation, and claim attribution.

\begin{longtable}[]{@{}
  >{\raggedright\arraybackslash}p{(\columnwidth - 4\tabcolsep) * \real{0.3333}}
  >{\raggedright\arraybackslash}p{(\columnwidth - 4\tabcolsep) * \real{0.3333}}
  >{\raggedright\arraybackslash}p{(\columnwidth - 4\tabcolsep) * \real{0.3333}}@{}}
\toprule\noalign{}
\begin{minipage}[b]{\linewidth}\raggedright
Dimension
\end{minipage} & \begin{minipage}[b]{\linewidth}\raggedright
Traditional review summarization / ABSA
\end{minipage} & \begin{minipage}[b]{\linewidth}\raggedright
Proposed system
\end{minipage} \\
\midrule\noalign{}
\endhead
\bottomrule\noalign{}
\endlastfoot
Output objective & Sentiment labels, topics, or free-form summaries & Structured VOC reports with evidence attribution \\
Product evidence access & Usually weak or absent & Joint access to text, images, and local visual regions \\
Review handling & Often aggregate statistics or direct summarization & Ratings are preserved and problem reviews are prioritized \\
Traceability & Error sources are difficult to locate & Intermediate artifacts, claim attribution, and feedback/replay are traceable \\
\end{longtable}

Table 2 gives an algorithmic description of the main workflow. The key point is that each stage emits persistent structured objects, so evaluation, replay, and error localization operate over the same object system.

\begin{longtable}[]{@{}
  >{\raggedright\arraybackslash}p{(\columnwidth - 6\tabcolsep) * \real{0.2500}}
  >{\raggedright\arraybackslash}p{(\columnwidth - 6\tabcolsep) * \real{0.2500}}
  >{\raggedright\arraybackslash}p{(\columnwidth - 6\tabcolsep) * \real{0.2500}}
  >{\raggedright\arraybackslash}p{(\columnwidth - 6\tabcolsep) * \real{0.2500}}@{}}
\toprule\noalign{}
\begin{minipage}[b]{\linewidth}\raggedright
Step
\end{minipage} & \begin{minipage}[b]{\linewidth}\raggedright
Input
\end{minipage} & \begin{minipage}[b]{\linewidth}\raggedright
Operation
\end{minipage} & \begin{minipage}[b]{\linewidth}\raggedright
Output
\end{minipage} \\
\midrule\noalign{}
\endhead
\bottomrule\noalign{}
\endlastfoot
1 & Raw product directory & Normalize product text, images, image regions, and reviews & Product evidence package \\
2 & Review records & Filter low-information reviews, perform rating-aware sampling, extract and deduplicate aspects & Review aspect objects \\
3 & Review aspect objects & Generate text-support, visual-confirmation, and cross-modal explanation questions & Route-aware retrieval questions \\
4 & Questions and evidence package & Retrieve candidates through text and image routes & Candidate evidence pool \\
5 & Candidate evidence pool & Select evidence by route, language, relevance, and reranking result & Ranked evidence bundle \\
6 & Aspect objects and evidence bundles & Aggregate aspect statistics and generate a structured report & Strengths, weaknesses, controversies, gaps, and suggestions \\
7 & Structured report and evidence objects & Perform claim-level attribution and evidence-gap labeling & Attributed report and run artifacts \\
\end{longtable}

\subsection{Product Evidence Representation}

Product pages contain names, descriptions, categories, specifications, images, and customer reviews. If these inputs are simply concatenated into a long prompt, later conclusions become difficult to trace. The system therefore builds product-level evidence packages in which text blocks, images, image regions, and reviews are represented as stable evidence objects.

Text evidence preserves source sections such as title, description, or category. Image evidence includes whole images and rule-based local regions, enabling local visual access without annotated bounding boxes. Review records preserve rating, language hints, and original text spans. The goal is not to complete semantic interpretation at the data layer, but to establish object boundaries for retrieval, attribution, and evaluation.

\subsection{Review Aspect Modeling}

The review-side task is to convert natural language feedback into aspect-level VOC units. Each aspect contains an attribute, sentiment, opinion, evidence span, usage scene, and confidence. Compared with whole-review summarization, aspect-level representation supports both statistical aggregation and evidence retrieval.

The system uses rating-aware sampling before extraction to avoid spending limited review budget on high-rating short comments. The default policy gives more weight to low-rating reviews because they often contain higher-value product improvement signals. Aspect extraction supports both structured LLM extraction and heuristic extraction, allowing the same workflow to run under full model settings or offline validation settings.

The default sampling profile in the current implementation is \texttt{problem\_first}. Given a maximum review budget, this profile assigns target weights of 30\%, 25\%, 20\%, 15\%, and 10\% to 1-star through 5-star reviews, respectively. When a rating bucket cannot satisfy its target quota, the remaining budget is filled according to the fallback order of 1-star, 2-star, 3-star, 4-star, and 5-star reviews. This design does not discard positive feedback; rather, it ensures that low-rating reviews are more likely to enter the aspect extraction stage when the analysis budget is limited, increasing the chance of detecting product defects, insufficient page presentation, and expectation mismatch.

Low-information comments are also handled around the sampling and extraction stage. The system preserves original reviews with ratings, but filters empty comments, overly short comments, and short low-information reviews such as \texttt{ok} and \texttt{good} before aspect extraction. This makes the review modeling layer depend not only on the number of reviews, but also on whether a review contains diagnostic information that can be converted into aspect objects and evidence-seeking questions. Each aspect entering later stages retains the attribute, sentiment, opinion, evidence span, usage scene, and confidence, providing a unified input for question planning and aspect aggregation.

\subsection{Aspect-to-Question Planning}

The key design is the intermediate planning layer between aspects and retrieval. Aspects such as small capacity, comfortable fit, or child appeal are still subjective and underspecified. Direct retrieval with these phrases often produces broad candidates and does not indicate whether text or images should be searched. The system therefore constructs evidence-seeking questions with explicit retrieval intents.

The questions cover text support, visual confirmation, and cross-modal resolution. Text-support questions ask whether product copy or specifications directly support a claim. Visual-confirmation questions ask whether product images show relevant structures or appearances. Cross-modal questions ask whether text and images together explain a review as a real product issue, insufficient presentation, or expectation mismatch.

\begin{figure}[H]
\centering
\includegraphics[width=0.92\textwidth]{/home/severin/Codelib/Decathlon\_VOC\_Analyzer/论文\_md形式/图片/图2.png}
\caption{Aspect-to-question planning. Aspect objects are not directly used as retrieval queries; they are decomposed into intent-specific questions that trigger text, image, or cross-modal evidence routes.}
\end{figure}

\subsection{Multimodal Recall and Reranking}

The evidence alignment layer uses two-stage retrieval. The first stage performs vector-based coarse recall over product text and product images. The second stage reranks the candidate set with higher-cost models. Text and image evidence are managed as separate routes during retrieval, but they are unified in the product evidence space during attribution.

Text evidence is useful for names, descriptions, specifications, and explicit claims. Image evidence is useful for structure, appearance, color, local details, and presentation quality. Whole images and rule-based local crops allow preliminary region-level evidence to enter the VOC pipeline. The candidate pool is balanced by route and language to prevent a single route or language from dominating final evidence selection.

To place textual relevance, visual relevance, and cross-modal consistency in a unified candidate-selection framework, the evidence score between a review aspect \(a\) and candidate evidence \(e\) can be abstracted as:

\[
S(a,e) = \lambda_t R_t(a,e) + \lambda_v R_v(a,e) + \lambda_c C(a,e)
\]

Here, \(R_t(a,e)\) denotes relevance on the text route, \(R_v(a,e)\) denotes relevance on the image route, and \(C(a,e)\) denotes the cross-modal consistency constraint between the review aspect, product text, and product images. The weights \(\lambda_t\), \(\lambda_v\), and \(\lambda_c\) are tunable. The formula does not require every implementation to use a single linear scoring model. Instead, it makes explicit the three evidence sources that candidate selection should consider: whether the product text makes an explicit commitment, whether the product image visibly supports the claim, and whether the two modalities jointly explain the customer experience. This reduces attribution errors such as treating text-only similarity as visual support or treating visually plausible evidence as a product promise.

\subsection{Evidence-Constrained Reporting and Claim Attribution}

The generation stage does not pass retrieved evidence to the model as ordinary context. It first aggregates aspect signals and retrieval coverage, then generates structured reports containing strengths, weaknesses, controversies, applicable scenes, evidence gaps, and suggestions.

To make the report auditable, the system performs claim-level attribution. Major claims are mapped to review evidence, product text evidence, product image evidence, or their combinations, and are labeled as supported, partially supported, unsupported, or contradicted. If a review claim cannot be verified by current product evidence, the system prefers to expose an evidence gap instead of presenting the claim as confirmed.

\begin{figure}[H]
\centering
\includegraphics[width=0.92\textwidth]{/home/severin/Codelib/Decathlon\_VOC\_Analyzer/论文\_md形式/图片/图3.png}
\caption{Evidence attribution and evaluation loop. Report claims are mapped to review, product text, and product image evidence, while evaluation measures both retrieval quality and claim grounding quality.}
\end{figure}

\subsection{Reproducibility and Explainability}

The reproducibility of the method comes from structured intermediate representations and persistent artifacts. Product evidence, review aspects, questions, retrieval records, report claims, and evaluation metrics are represented as typed objects. This allows researchers to freeze earlier stages and replace question generation, retrieval, or reranking strategies for ablation.

The same design improves error analysis. If a final suggestion is unreliable, the error can be traced to sampling, extraction, question planning, candidate recall, reranking, report generation, or attribution calibration. This makes the framework more suitable for human review and iterative system research than black-box summarization.

At the persisted artifact level, the system saves the analysis report, review extraction output, question intents, retrieval questions, retrieval records, retrieval quality, runtime configuration, aspect aggregates, process traces, and replay summaries. The manifest records run configuration, stage artifacts, and evaluation summaries. The feedback sidecar organizes weaknesses, suggestions, and retrieval-quality issues into slots for human review. The replay sidecar stores the previous report, process trace, and retrieval quality so that later runs can identify persistent issues, resolved issues, and new issues, and can reflect accepted or rejected human feedback in report ordering and suggestion notes.

The system also applies signature-bound caching to query embeddings and rerank results. A query embedding cache signature includes the retrieval route, query text, backend type, model name, and service endpoint. A rerank cache signature additionally includes candidate count, candidate digest, reranking backend, and model configuration. The candidate digest is computed from fields such as evidence ID, product ID, route, text block, image, region, language, and score. These signature constraints reduce repeated computation while preventing cache contamination across different backends, models, or candidate sets.

% =========================
% Experimental Setup
% =========================
\section{Experimental Setup}
\subsection{Experimental Goal}

The experiments validate whether the proposed system can perform evidence-driven VOC analysis reliably. We evaluate whether the system can generate complete structured reports from raw product data, whether intermediate objects support audit and error localization, whether full model paths and degraded paths are distinguishable, and whether the evaluation interface can record retrieval quality and report attribution quality.

This setup matches the system-methodology nature of the paper. The paper establishes a reproducible protocol in which the same input product, run configuration, intermediate artifacts, and evaluation summaries can be saved, inspected, and replayed.

\subsection{Implementation Environment}

The system is implemented in the Python ecosystem with both Web API and batch workflow interfaces. Typed schemas constrain intermediate objects, a state graph organizes the multi-stage workflow, and language model calls are routed through a compatible model gateway. The retrieval layer supports local and vector-database backends, image evidence can be encoded with vision-language representations, and candidate evidence can be reranked by text or multimodal rerankers.

Runtime policies distinguish formal evaluation from development validation. When full model dependencies are available, the system uses real embeddings, image encoders, and rerankers. When external capabilities are unavailable and degradation is allowed, the system can fall back to heuristic paths to keep the workflow executable.

The system exposes both interactive API and offline script entry points. The API layer provides endpoints for dataset overview, product normalization, index overview, index construction, review aspect extraction, and single-product analysis. Batch scripts execute the full workflow, offline validation, multimodal run checks, HTML export, manifest writing, and experiment matrices. Because these entry points share the same schemas and service layer, researchers can inspect individual stages through the API during development and reproduce the same workflow in batch mode during experiments.

Table 3 summarizes implementation details directly related to reproducible experiments.

\begin{longtable}[]{@{}
  >{\raggedright\arraybackslash}p{(\columnwidth - 2\tabcolsep) * \real{0.5000}}
  >{\raggedright\arraybackslash}p{(\columnwidth - 2\tabcolsep) * \real{0.5000}}@{}}
\toprule\noalign{}
\begin{minipage}[b]{\linewidth}\raggedright
Item
\end{minipage} & \begin{minipage}[b]{\linewidth}\raggedright
Current implementation
\end{minipage} \\
\midrule\noalign{}
\endhead
\bottomrule\noalign{}
\endlastfoot
Language & Python 3.11 or above \\
Workflow entry points & Web API and \texttt{run\_workflow.py} batch script \\
Workflow orchestration & LangGraph state graph \\
Data object constraints & Pydantic schemas \\
Retrieval backends & Local JSON index and Qdrant backend \\
Text and image representations & Text embeddings, CLIP or compatible vision-language embeddings, heuristic fallback paths \\
Reranking paths & Text reranker, multimodal reranker, or heuristic ranking \\
Cached objects & Query embeddings, rerank results, and analysis checkpoints \\
Main artifacts & Normalized evidence package, aspect objects, retrieval records, structured reports, feedback/replay sidecars, HTML, and manifests \\
Engineering validation & Tests for normalization, review modeling, retrieval, report generation, manifest evaluation, and workflow entry points \\
\end{longtable}

\subsection{Data and Validation Units}

The data consists of product-page crawl results, including product text, product images, and customer reviews. The system first converts raw inputs into product evidence packages, then executes single-product analysis. The validation unit is therefore a single product run rather than an aggregate cross-product metric. Each run produces structured records for aspects, questions, retrievals, reports, and evaluation.

The current validation unit does not support cross-product performance claims. Its role is to confirm that the system can produce complete, auditable, and replayable analysis samples at single-product granularity.

\subsection{Metrics}

The evaluation metrics are divided into two groups. Retrieval metrics include Recall@1, Recall@3, Recall@5, MRR, NDCG@3, and NDCG@5 when question-level evidence labels are available. These are standard ranking and recall metrics in information retrieval \citep{manning2008introductionIr}.

Claim-grounding metrics evaluate whether structured report claims are supported by review, product text, or product image evidence. They include claim support rate, claim grounded rate, citation precision, citation contradiction rate, and modality contribution. These metrics ask whether the report is evidence-supported rather than merely fluent.

\subsection{Engineering Validation Protocol}

Automated testing is part of the validation protocol. The tests cover data normalization, review modeling, question generation, index backends, embedding and reranking, retrieval, report generation, HTML export, manifest evaluation, runtime policies, and workflow entry points. The current codebase contains 166 passing tests, indicating that the main interfaces, structured artifacts, and workflow assertions are aligned.

% =========================
% Results and Analysis
% =========================
\section{Results and Analysis}
\subsection{Workflow Completeness}

The first validation result is that the current implementation completes the full analysis chain from raw product information to attributed structured reports. It performs product evidence normalization, review aspect modeling, question planning, multimodal recall, candidate reranking, aspect aggregation, report generation, and evidence attribution. Each stage emits consistent structured objects, so a single run can be decomposed into auditable analytical units.

This demonstrates that the method is not only conceptual. Review-side objects, product-side evidence, and report-side claims are connected through explicit data flow. Even when external model calls are disabled for validation, the system preserves the same object boundaries, separating workflow correctness from model capability.

\subsection{Structured Auditability}

A complete analysis contains not only a final narrative report but also review aspects, question intents, retrieval questions, retrieved evidence, retrieval quality, runtime configuration, aspect aggregation, process trace, and report attribution. These objects allow analysts to inspect the source of a conclusion instead of accepting a single opaque summary.

This structure directly supports error localization. If a suggestion lacks support, one can inspect whether the review was extracted correctly, whether the question was too broad, whether text or image retrieval found valid evidence, whether reranking prioritized key evidence, and whether report refinement correctly identified evidence gaps.

\subsection{Role of Question Planning}

Question planning is the key structure connecting reviews and product evidence. Raw reviews often contain emotion, context, and implicit assumptions; direct retrieval can scatter evidence candidates. Aspect-level questions convert these expressions into clearer evidence needs, such as whether product copy explicitly supports an experience, whether images show the relevant structure, or whether cross-modal evidence explains a negative review.

Because every retrieval keeps its source aspect, source question, and expected evidence routes, the system can be analyzed at query level. Questions, text blocks, images, regions, and report claims share stable identifiers, so analysts can inspect whether a question is supported by evidence from the expected route.

\subsection{Complementarity of Multimodal Evidence}

Text and image evidence serve different functions in product VOC analysis. Text evidence is suited for explicit specifications, usage claims, categories, and descriptions. Image evidence is suited for structure, appearance, color, local details, and presentation quality. The route-aware design allows the report to state whether a review claim is supported by text, by images, by both, or by neither.

This design also prevents image evidence from being textified too early. For appearance and structure-related issues, image candidates can be reranked by a vision-language model that directly inspects the whole image or crop. Although the current regions are rule-based, they demonstrate the feasibility of introducing local visual evidence into VOC analysis.

\subsection{Extensibility of Evaluation}

The evaluator supports both retrieval quality and report-grounding quality. On the retrieval side, it records questions, candidate evidence, ranking scores, and route coverage. On the report side, it records claim support status, cited evidence, contradiction risk, and modality contribution. System validation therefore checks not only whether final text is generated, but also whether conclusions retain traceable evidence.

The evaluation interface shares identifiers with intermediate objects. Questions, text blocks, images, regions, and report claims can enter review, replay, and error analysis without being re-encoded.

\subsection{Boundary of Current Results}

The current results establish workflow completeness, artifact auditability, and a stable test baseline. This section does not provide cross-product performance rankings. It limits the conclusion to system-method validation: whether the system can reliably generate traceable artifacts and connect review aspects, retrieval questions, product evidence, and report claims in one analysis chain.

% =========================
% Discussion
% =========================
\section{Discussion}
\subsection{Why Question-Driven Retrieval Fits Product VOC}

The difficulty in product VOC analysis is not only summarizing reviews, but also explaining whether review claims are supported by product evidence. Question-driven retrieval adds an intermediate semantic layer: review aspects describe customer concerns, while retrieval questions describe evidence needs that product text or images can answer.

This is suitable for e-commerce because product pages are naturally multimodal and reviews are naturally subjective and noisy. Directly concatenating all information into a model sacrifices controllability, whereas question planning turns feedback into retrievable, evaluable, and attributable units.

\subsection{Value of Artifact-First Engineering}

The system is artifact-first. Normalized evidence packages, aspect extraction results, question caches, retrieval caches, reports, feedback, replay files, HTML exports, and manifests have fixed structures and output locations. Researchers can inspect any stage instead of only reading the final report.

This also lowers review cost. Standardized product packages, aspect objects, retrieval questions, HTML reports, and manifests can be reused to inspect the same product run stage by stage without rerunning the complete model chain.

\subsection{Comparison with One-Step LLM Summarization}

One-step LLM summarization is simple, but it cannot reliably answer where a conclusion came from, whether product text supports it, or whether images support it. Our staged design increases implementation complexity but provides interpretability, reproducibility, and evaluability.

For a research-oriented system, this tradeoff is necessary. Product VOC reports must support evidence inspection if they are to be used in operational or product decisions.

\subsection{Methodological Implications}

The implementation suggests that better evidence organization and query planning may be more important than simply using larger generators. Review aspects, evidence-seeking questions, text routes, image routes, and claim attribution form auditable operation units, allowing the system to organize product VOC analysis without relying on one black-box summary call.

\subsection{Scope of Applicability}

The method is suitable for product analysis scenarios where evidence provenance matters, especially when product pages contain text, images, and reviews. For scenarios that only require quick high-level summaries, the multi-stage system may be heavier than necessary. For highly localized visual quality issues, the current whole-image and rule-based region evidence still requires finer-grained visual models.

% =========================
% Conclusion
% =========================
\section{Conclusion}
This paper presents Decathlon VOC Analyzer, an evidence-driven multimodal VOC analysis system for aligning product text and image evidence with customer reviews. The system normalizes raw product data into traceable evidence packages, extracts aspect-level VOC objects from reviews, plans evidence-seeking questions, retrieves and reranks product text and image evidence, and generates structured reports with claim-level attribution.

The central conclusion is that product VOC analysis should not be reduced to one-step review summarization. It is better organized as a multi-stage workflow: review aspect modeling, evidence question planning, product text-image retrieval, and evidence-constrained generation. The current implementation forms a runnable research prototype with consistent structured intermediate representations, evidence attribution, replay mechanisms, and evaluation interfaces. The current codebase passes 166 automated tests.

The system also preserves original reviews with ratings and strengthens the role of low-rating feedback through problem-first rating-aware sampling and low-information review filtering. This design helps the system represent product-improvement signals more reliably beyond general review summarization, and provides a clearer evidence entry point for distinguishing product defects, missing page information, and expectation mismatch.

Overall, the paper shows a feasible path for turning product VOC analysis from black-box summarization into an evidence workflow. The key contribution is not a single generated report, but the structured connection among reviews, questions, evidence, reports, and evaluation.

% =========================
% Limitations
% =========================
\section{Limitations}
First, the current results are system validation results. The system supports manifest evaluation and automated testing, but this paper does not build a frozen multi-category human-labeled dataset and therefore does not report cross-product performance rankings.

Second, this paper does not report retrieval-strategy ablations. The implementation supports text routes, image routes, question-driven retrieval, reranking, and caching, but these capabilities are described as system components rather than ranked against one another.

Third, visual evidence granularity is limited. The system supports whole images and fixed-ratio crops, but these regions are not produced by semantic detection, segmentation, or grounding models.

Fourth, LLM outputs still depend on prompts and schema constraints. Structured output, refinement, and heuristic fallbacks reduce but do not eliminate misunderstanding, omission, or overgeneralization.

Fifth, the feedback loop is still mostly an engineering interface. Feedback and replay sidecars exist, but this paper does not quantify the quality impact of replay on later reports.

Sixth, image evidence is still based on whole images and rule-based regions. The relevant claims should be understood as a system interface for local visual evidence rather than proof of precise visual grounding.

% =========================
% Acknowledgments
% =========================
\section{Acknowledgments}
This paper is based on the source implementation, design documents, test scripts, and run artifacts of the Decathlon VOC Analyzer project. We thank the project members who provided feedback and assistance during development and validation.

% =========================
% References
% =========================
\bibliographystyle{unsrtnat}
\bibliography{ref}

\end{document}
